\scanpage{57}
Comme on a supposé 2 régulier dans $S$, ceci \deleted{\unsure{impliq|implique}} détermine $\varpi$ de façon unique (d'où \uline{unicité} \deleted{de} $q$). Pour la condition d'existence, posons
\[S_1=\unsure{\mathrm{Var}|\mathrm{Van}|V}(4\underline O_S)\subset S\tag*{(112)}\]
(premier voisinage infinitésimal de $S_0$), et considérons. $\Theta(q_0)-D\,|\,S_1=a_1\deleted{\unsure{(q)}}$, il ne dépend pas du choix de $q_0$, et est de lui \unsure{globalt|globalemt|globale} canonique [défini modulo l'existence -- et choix, le cas échéant -- de $b_0$ satisfaisant (106)]
\[a_1(D)=a_1\in\Gamma(S_1,L^{\otimes2}),\tag*{(113)}\]
et la condition \deleted{est que} d'existence de $q$ est
\[a_1(D)=0\tag*{(114)}\]

On a prouvé :

\uline{Th} Soit $S$ schéma tel que $2\cdot1_S$ soit $S$-régulier, $S_0=V(2\cdot\underline O_S)$, $S_1=V(4\underline O_S)$ son premier voisinage infinitésimal. \deleted{On se donne un \unsure{M.} inv. \unsure{soit,} une section $D$ de $L$} \inserted{Alors la cat. des \deleted{ext.} \inserted{rev.} quadratiques} de $S$ (ou des ext. quadratiques de $\mathcal O_S$) équivaut à celle des triples $(L,D,b_0)$ d'un module inversible $L$, \deleted{et} d'une section $D$ de $L^{\otimes2}$ \deleted{satisfaisant} \inserted{et de $b_0\in\Gamma(S_0,L|S_0)$, satisfaisant} les 2 conditions suivantes

a) \deleted{$\exists\,b_0\in\Gamma(S_0,L|S_0)$ telle que} $b_0^2=D_0\overset{\text{déf}}=D|S_0$

b) La section $a_1(D)\in\Gamma(S_1,L^{\otimes2}|S_1)$ (définie plus haut \deleted{est nulle} en termes de $D_1=D|S_1$) est \uline{nulle}.
% archive sheet number: 56
